\chapter{Consumo energetico}
\label{cap:consumi}

\begin{quote}\small
Questo capitolo risponde all'obiettivo O4. L'analisi è dichiaratamente
\emph{a titolo informativo} e basata su dati di datasheet: non si dispone di
strumentazione di misura (né multimetro né resistenze di shunt), come concordato
con il relatore. Tutti i valori riportati sono quelli dichiarati dai produttori,
con la fonte indicata.
\end{quote}

\section{Consumo dei singoli componenti}

\subsection{Sensori}

\begin{table*}[htbp]
\centering
\caption{Consumo dei sensori in esame, da datasheet.}
\label{tab:consumi-sensori}
\small
\begin{tabularx}{\linewidth}{@{}lllrX@{}}
\toprule
\textbf{Componente} & \textbf{Tensione} & \textbf{Corrente media} &
\textbf{Potenza} & \textbf{Fonte} \\
\midrule
\ldb & 5~V ($>$200~mA di capacità) & 80~mA & $\sim$400~mW &
  manuale Hi-Link V1.03~\cite{hilinkManualV103} \\
\ldc & 3,3~V (3,0--3,6~V) & 50~mA & $\sim$165~mW &
  manuale Hi-Link V1.2~\cite[Tab.~2-1]{hilinkManualLD2420} \\
PIR \pir & 4,5--20~V & $<$0,05~mA & $\sim$0,3~mW &
  datasheet \pir~\cite{hcsr501} \\
\bottomrule
\end{tabularx}
\end{table*}

Il rapporto fra i due estremi è immediato e va sottolineato perché è il dato
strutturale di tutto il capitolo: il radar consuma circa \textbf{1600 volte} la
corrente del PIR.

\paragraph{Nota su una fonte discordante.} Una fonte secondaria di uso
comune~\cite{components101} riporta per l'\pir\ una corrente di riposo di 65~mA. Si tratta con evidenza di un
refuso fra milliampere e microampere, incompatibile con l'architettura del modulo
(un BISS0001 alimentato da un LDO HT7133). Si è fatto fede al datasheet, che indica
$<$50~\textmu A. Il caso è riportato perché illustra la necessità della regola
adottata in tutta la tesi: risalire sempre al datasheet del produttore.

\subsection{Microcontrollore}

\begin{table*}[htbp]
\centering
\caption{Consumo dell'ESP32-WROOM-32 per modalità operativa
(datasheet Espressif~\cite{esp32datasheet}).}
\label{tab:consumi-esp32}
\small
\begin{tabularx}{\linewidth}{@{}llX@{}}
\toprule
\textbf{Modalità} & \textbf{Corrente} & \textbf{Note} \\
\midrule
Attivo, WiFi in trasmissione & 160--260~mA (picco) & invio dati alla web UI \\
Attivo, WiFi in ricezione/idle & $\sim$95--100~mA & connesso, in ascolto \\
Modem-sleep (CPU attiva, WiFi spento) & 20--68~mA & dipende dalla frequenza di
  clock (80--240~MHz) \\
Light-sleep & $\sim$0,8~mA & \\
Deep-sleep & $\sim$0,01~mA (10~\textmu A) & solo dominio RTC attivo \\
\bottomrule
\end{tabularx}
\end{table*}

\section{Consumo dei nodi completi}

Sommando microcontrollore e sensore si ottengono le configurazioni effettivamente
usate in questa tesi e quelle rilevanti per il progetto (\cref{tab:consumi-nodi}).

\begin{table*}[htbp]
\centering
\caption{Consumo stimato del nodo completo nelle configurazioni di interesse.}
\label{tab:consumi-nodi}
\small
\begin{tabular}{@{}lrr@{}}
\toprule
\textbf{Configurazione} & \textbf{Corrente totale} & \textbf{Potenza (5~V)} \\
\midrule
ESP32 attivo + LD2410B + WiFi & $\sim$180--260~mA & $\sim$0,9--1,3~W \\
ESP32 attivo + LD2410B (solo logging seriale) & $\sim$110--150~mA & $\sim$0,55--0,75~W \\
ESP32 attivo + LD2420 & $\sim$80--120~mA & $\sim$0,4--0,6~W \\
ESP32 attivo + PIR & $\sim$30--70~mA & $\sim$0,15--0,35~W \\
ESP32 in deep-sleep + PIR (sveglia su interrupt) & $\mathbf{\sim 0{,}06}$~\textbf{mA} & $\sim$0,3~mW \\
\bottomrule
\end{tabular}
\end{table*}

L'ultima riga è il punto chiave del capitolo. L'uscita digitale del PIR può
risvegliare l'ESP32 dal deep-sleep tramite interrupt su GPIO: il nodo può quindi
restare in attesa a costo praticamente nullo, con il PIR nel ruolo di
\emph{guardiano}. Il radar mmWave, invece, costa tre ordini di grandezza in più ed
è progettato per un'alimentazione fissa.

\section{Stime di autonomia a batteria}

Assumendo una cella litio-ione 18650 da 3000~mAh e un rendimento del regolatore di
circa l'85\,\%, si ottengono le autonomie di \cref{tab:autonomia}.

\begin{table*}[htbp]
\centering
\caption{Autonomia stimata con una cella Li-ion 18650 da 3000~mAh e rendimento del
regolatore dell'85\,\%.}
\label{tab:autonomia}
\small
\begin{tabular}{@{}lrr@{}}
\toprule
\textbf{Scenario} & \textbf{Corrente media} & \textbf{Autonomia} \\
\midrule
Nodo mmWave sempre attivo (LD2410B + ESP32 + WiFi) & $\sim$200~mA & $\sim$13~h \\
Nodo mmWave attivo senza WiFi continuo & $\sim$130~mA & $\sim$20~h \\
Nodo PIR sempre attivo (ESP32 in modem-sleep) & $\sim$25~mA & $\sim$4 giorni \\
Nodo dormiente (deep-sleep + PIR come watchdog) & $\sim$0,06~mA & $\sim$4--5 anni$^{*}$ \\
\bottomrule
\multicolumn{3}{@{}l@{}}{\footnotesize $^{*}$ in questo regime il fattore
limitante non è il consumo ma l'autoscarica della cella.}
\end{tabular}
\end{table*}

\section{Implicazione per UPRISE: l'architettura ibrida}
\label{sec:architettura-ibrida}

I numeri precedenti risolvono un'apparente contraddizione del progetto. Il sistema
integrato negli arredi deve restare in attesa \emph{per anni} in ``tempo di pace'' e
deve funzionare a batteria durante il blackout che segue il sisma: due requisiti
incompatibili con un sensore da 80~mA sempre acceso, ma anche con un sensore che non
vede le persone ferme.

La soluzione non è scegliere fra i due sensori, ma metterli in sequenza temporale.

\begin{enumerate}[leftmargin=1.8em,itemsep=4pt]
  \item \textbf{Tempo di pace} --- l'intero nodo in deep-sleep, con il PIR spento o
        usato come watchdog: consumo dell'ordine dei microampere, autonomia
        pluriennale limitata dall'autoscarica della batteria.
  \item \textbf{Trigger di ``modalità terremoto''} --- l'accelerometro sul gateway
        o il rilevamento del blackout risveglia l'ESP32. Il PIR può contribuire come
        sorgente di interrupt a costo nullo.
  \item \textbf{Emergenza} --- si accende il radar mmWave per il rilevamento fine:
        persona ferma sotto il banco, respiro, indice di vitalità. L'autonomia di
        13--20~ore copre la finestra critica iniziale delle operazioni di ricerca e
        soccorso.
  \item \textbf{Estensione della finestra} --- per coprire le prime 72~ore, che è
        l'orizzonte convenzionale del soccorso, sono disponibili due strade
        combinabili: una batteria maggiorata, oppure il \emph{duty-cycling} del
        radar. Un ciclo di 1~minuto acceso ogni 5 riduce il consumo medio di circa
        cinque volte, portando l'autonomia oltre le 72~ore. Il costo è una latenza
        di rilevamento fino a quattro minuti, che nello scenario di soccorso è
        accettabile: la domanda a cui il sistema risponde è ``c'è qualcuno sotto
        questo banco?'', non ``in quale istante si è mosso?''.
\end{enumerate}

\paragraph{Conclusione del capitolo.} L'analisi dei consumi porta a una
raccomandazione precisa, che è anche la risposta alla domanda implicita del titolo
della tesi: il PIR \textbf{non va sostituito}. Non è un concorrente del radar ma il
guardiano a basso costo che decide quando accenderlo. Il radar mmWave è lo strumento
di misura della fase di emergenza. Ciascuno dei due fa ciò che l'altro non può fare,
e i dati del \cref{cap:confronto} e di questo capitolo lo mostrano da due direzioni
opposte.

\todo{se il tempo lo consente: verificare sperimentalmente il consumo con un
misuratore USB in linea, anche a bassa precisione. Anche una misura al 10\,\%
di accuratezza confermerebbe l'ordine di grandezza delle tabelle e renderebbe il
capitolo più solido}
